\def\HostName{Codex}
\documentclass[UTF8,a4paper,11pt,fontset=fandol]{ctexart}
\usepackage[margin=22mm,headheight=14pt]{geometry}
\usepackage{amsmath,amssymb,booktabs,longtable,array,xcolor,fvextra,fancyhdr,enumitem}
\usepackage[colorlinks=true,linkcolor=blue!55!black,urlcolor=blue!60!black,breaklinks=true]{hyperref}
\setmonofont{Latin Modern Mono}
\setCJKmonofont{FandolFang-Regular}
\definecolor{ink}{HTML}{15324F}
\definecolor{light}{HTML}{EEF3F8}
\DefineVerbatimEnvironment{code}{Verbatim}{fontsize=\small,breaklines=true,breakanywhere=true,frame=single,rulecolor=\color{ink!25},framesep=3mm}
\setlist{itemsep=2pt,topsep=4pt}
\setlength{\parindent}{2em}
\setlength{\parskip}{4pt}
\setlength{\emergencystretch}{3em}
\pagestyle{fancy}\fancyhf{}
\fancyhead[L]{\small CryptoBoolFunc 第一版}
\fancyhead[R]{\small\HostName 教程}
\fancyfoot[C]{\thepage}
\hypersetup{pdftitle={CryptoBoolFunc 第一版安装与研究教程},pdfauthor={wernerchang; Codex (AI collaborator)}}
\newcommand{\note}[1]{\par\noindent\colorbox{light}{\parbox{0.94\linewidth}{\small #1}}\par}
\newcommand{\file}[1]{\texttt{\small\detokenize{#1}}}
\title{\color{ink}CryptoBoolFunc 第一版\\\large 从安装仓颉到函数构造的保姆级教程\\\HostName 安装版本}
\author{项目维护：wernerchang\\AI 协作贡献：Codex（OpenAI）}
\date{版本 1.0.0（研究试用）\quad 2026年11月5日\\\small 资料核查：2026年11月4日}
\begin{document}
\sloppy\raggedright
\maketitle
\note{操作主线：先安装完整仓颉 cangjie-skill，再安装本项目第一版；然后确认宿主能读文件，运行一个数值实例，最后进入研究任务。两种安装版本共用相同数学内容。}
\tableofcontents
\clearpage

\section{第一步：先把仓颉安装材料准备好}
\subsection{先分清两个名称}
仓颉 \file{cangjie-skill} 是上游“把长内容蒸馏成技能”的方法与工具。CryptoBoolFunc 是利用该流程整理出的密码函数研究技能。前者负责提取、验证、编译；后者包含本书的15张能力卡和小参数计算程序。只下载本项目，并不等于已经安装了仓颉。

本教程以 Windows PowerShell 为主。先打开开始菜单，搜索 PowerShell，进入普通终端。命令按块执行，每块结束先看“成功标志”，不要一次粘贴整篇。macOS/Linux 的差异在附录列出。所有下载采用实际发布附件；不要把网页另存的HTML改名为ZIP。

\subsection{检查 Python 与工作目录}
需要 Python 3.10 或以上和网络。运行：
\begin{code}
python --version
python -m pip --version
\end{code}
成功标志是显示 Python 版本及pip路径。如果打开微软商店、提示找不到命令，先从\href{https://www.python.org/downloads/}{Python官网}安装，勾选添加PATH，再关闭并重新打开终端。如果电脑只有 \file{py -3} 可用，本文所有 \file{python} 统一换成 \file{py -3}，不要混用不同解释器安装依赖。

创建实验目录，后面的文件都放在这里：
\begin{code}
$lab = Join-Path $env:USERPROFILE 'CryptoBoolFunc-Lab'
New-Item -ItemType Directory -Path $lab -Force | Out-Null
Set-Location $lab
Get-Location
\end{code}
成功标志是当前路径末尾为 \file{CryptoBoolFunc-Lab}。它是示例工作目录，可改到自己的数据盘；路径含空格时保留引号。

\subsection{下载、验哈希、解压完整仓颉}
本版教程采用上游v2.5.0在2026-09-13刷新后的通用附件。上游说明，同一标签自动生成的Source code归档与刷新附件可能不同。以附件和校验值为准\cite{cangjie}。

\begin{code}
$base = 'https://github.com/kangarooking/cangjie-skill/releases/download/v2.5.0'
Invoke-WebRequest "$base/cangjie-skill-2.5.0.zip" -OutFile cangjie.zip
Invoke-WebRequest "$base/cangjie-skill-2.5.0.zip.sha256" -OutFile cangjie.sha256
$actual = (Get-FileHash cangjie.zip -Algorithm SHA256).Hash.ToLower()
$expected = ((Get-Content cangjie.sha256 -Raw).Trim() -split '\s+')[0].ToLower()
if ($actual -ne $expected) { throw 'SHA256不匹配，请重新下载' }
Expand-Archive -LiteralPath cangjie.zip -DestinationPath upstream
$cj = Join-Path $lab 'upstream\cangjie-skill'
Test-Path (Join-Path $cj 'SKILL.md')
Test-Path (Join-Path $cj 'scripts\cangjie.py')
\end{code}
两项都应为 \file{True}。若目录已存在，先换一个新的解压目录，不要直接覆盖旧研究文件。已核实附件SHA256为：
\begin{code}
c08bd78239fde0593dffdd5c315f9eceb59c22872baa5df54d34798449070906
\end{code}
如果上游以后重新发布而附件变化，不应忽略不匹配；应记录新的版本与来源，重新验证。

\subsection{安装依赖并自检}
\begin{code}
python -m pip install pyyaml
python "$cj\scripts\cangjie.py" doctor
\end{code}
成功标志是 \file{doctor: PASS}。\file{tiktoken}、\file{jsonschema} 为可选项；它们的缺失提示与PyYAML缺失不同。若报 \file{No module named yaml}，用同一个 \file{python -m pip} 重装后再测。此时只是仓颉文件和工具就绪；下一节才将完整目录放入对应宿主扫描根。

\section{第二步：将仓颉安装进 Codex}
\subsection{准备宿主与登录}
如果已有Codex桌面应用，打开并确认能正常开始本地任务。如果尚未安装，从\href{https://developers.openai.com/codex/quickstart}{官方快速开始}进入当前系统对应的安装/登录流程\cite{codexquick}。本教程不要求切换模型或修改现有账号设置。

根据核查日官方指南，Codex用户级本地技能目录为 \file{~/.agents/skills}，项目级为仓库内 \file{.agents/skills}；应安装完整目录\cite{codexskills}。旧机器可能有兼容目录，不要把同名技能在多个根重复安装，以免选择器出现重复项。

\subsection{安装完整仓颉目录}
接着在刚才PowerShell中运行：
\begin{code}
$skillRoot = Join-Path $env:USERPROFILE '.agents\skills'
New-Item -ItemType Directory -Path $skillRoot -Force | Out-Null
$dst = Join-Path $skillRoot 'cangjie-skill'
if (Test-Path $dst) { throw '已存在仓颉，请先阅读升级与备份章节' }
Copy-Item -LiteralPath $cj -Destination $dst -Recurse
Test-Path (Join-Path $dst 'methodology')
Test-Path (Join-Path $dst 'scripts\cangjie.py')
\end{code}
成功标志是两项True，且 \file{cangjie-skill/SKILL.md} 只有一层技能目录。错误结构是 \file{cangjie-skill/cangjie-skill/SKILL.md}。只复制一个SKILL.md会导致后续模板和编译脚本缺失。

\subsection{确认 Codex 真的发现并读取}
新开一个Codex本地会话，输入：
\begin{code}
请使用 $cangjie-skill。
先读取该技能，告诉我安装路径、版本与doctor检查方法。
此时不要开始蒸馏，也不要修改文件。
\end{code}
成功标准是能指出实际安装目录并找到方法文件；只说“我会帮你”不算发现成功。CLI/IDE可使用技能选择器或显式名称；界面变化时以当前官方说明为准。未出现时先重开会话或重启Codex，再检查目录、frontmatter和权限。

\section{第三步：安装 CryptoBoolFunc 第一版}
\subsection{获取仓库与两种下载方式}
推荐安装Git后克隆：
\begin{code}
Set-Location $lab
git clone https://github.com/wernerchang/CryptoBoolFunc.git
Set-Location CryptoBoolFunc
Get-Content VERSION
\end{code}
第一次发布应显示 \file{1.0.0}。如果以后main已更新，请从仓库版本目录下载第一版Codex ZIP，并对照SHA256SUMS；不要用后来main的内容冒充第一版。不会使用Git时，可在仓库下载ZIP并解压到实验目录，确认根目录有 \file{scripts/install.py} 和 \file{skills}。

\subsection{先预演，再实际安装}
\begin{code}
python scripts/install.py --host codex --scope user --dry-run
python scripts/install.py --host codex --scope user
\end{code}
预演输出 \file{dry_run}，不复制文件。实际安装输出 \file{installed}；内容相同再次执行会输出 \file{already_installed}。若现有目录内容不同，会拒绝覆盖；确认升级后按后文备份方式执行。

项目级安装则用：
\begin{code}
python scripts/install.py --host codex --scope project --project .
\end{code}
用户级与项目级选一种即可。项目级的好处是只影响当前仓库；用户级适用于多个研究目录。第一版Codex ZIP中包含相同的安装脚本和完整技能目录；它不是需要在Python中 \file{pip install} 的库。

\subsection{新会话验收}
在Codex中打开克隆的仓库作为工作区，新开会话并输入：
\begin{code}
请使用 $mmhs-cryptographic-function-design。
先读取入口，列出它的15张能力卡；说明DS只允许R_z=I，
以及同偏差校正阶和弹性阶的区别。
请给出实际读取路径，不要开始构造。
\end{code}
成功标准：找到入口、15卡以及明确边界。若回答只有目录名称而不读取，明确要求打开SKILL.md及相应卡。模型能口头复述不等于程序已执行；下一章用实际数值验收。

\section{第四步：先运行一个能核对的数值实例}
\subsection{MM：验证本机计算链路}
以下命令在CryptoBoolFunc仓库根执行。它不调用收费模型API，也不需要把原书上传到网络：
\begin{code}
New-Item -ItemType Directory -Path output -Force | Out-Null
$tool = 'skills/mmhs-cryptographic-function-design/resources/math_tool.py'
python $tool mm --input examples/mm.json --output output/mm-result.json
$r = Get-Content output/mm-result.json -Raw | ConvertFrom-Json
$r.status
$r.metrics.nonlinearity
$r.metrics.balanced
$r.metrics.sac
\end{code}
应依次看到ok、4、True及4个True。样例的输入为k=3、选择器系数3与4、常数项1与0，得到4变量函数。不要将输入选择器长度2误解为2位选择器。工具采用整数bit i表示第i输入/输出坐标，原书其他排表方式须先换坐标再对比。

\subsection{多输出：必须检查所有非零组合}
\begin{code}
python $tool ds --input examples/ds.json --output output/ds-result.json
$r = Get-Content output/ds-result.json -Raw | ConvertFrom-Json
$r.metrics.component_count
$r.metrics.min_nonlinearity
$r.metrics.balanced
\end{code}
应看到7、240、True。因为m=3，需要检查7个非零输出组合。\file{min_nonlinearity} 是最坏分量非线性度，\file{min_degree} 与 \file{max_degree} 分别报告两个次数口径。本版DS只实现经过复核的 \(R_z=I\) 子族；指定自定义R会拒绝。一般可逆R分支已发现反例，不能因为示例跑通而恢复其一般结论。

\subsection{校正器与阵列例}
\begin{code}
python $tool corrector --input examples/corrector.json --output output/corrector.json
python $tool d-array --input examples/d-array.json --output output/d-array.json
\end{code}
双码校正器应报告最小非线性度448、最小次数4、整体校正阶4。它要求输入位统计独立且具有相同偏差；若模型条件改变，不能沿用相同保证。阵列例固定多项式37，即 \(x^5+x^2+1\)，合法行数应为10；换为已核查的另一表示47时可为11。容量要连同多项式、基和列选择一起记录。

\section{第五步：让宿主按卡执行，而非只口头解释}
在已打开正确工作区的宿主会话中提交：
\begin{code}
请使用mmhs-cryptographic-function-design中的mm-spectrum卡。
输入取examples/mm.json。
先列合法性条件，再运行资源程序，把输出写到output/mm-session.json。
报告实际命令、退出码、完整表、Walsh谱、非线性度与SAC。
区分原书结论、本文重建和本次有限实例；不要扩大为一般最优。
\end{code}
成功标准有四项：读取正确卡；实际运行命令；输出文件存在且可解析；数值与前面手工验收一致。宿主如果需要权限批准，应核对具体命令和输出路径。API失败、文件权限失败与数学条件不满足是三种不同失败，应分别记录。

再做一次边界验收：把MM输入的phi改成重复值，例如[1,1]，使用另一个输入文件运行。工具应返回error且退出码2，不能继续报告单射构造成功。无论哪个宿主，“明确拒绝非法输入”都是必要行为。

\section{第六步：从新材料重新使用仓颉蒸馏}
如果只是使用本版15项方法，不必重新蒸馏。若要扩展到新论文、另一份书稿或新章节，将合法持有的PDF留在本机 \file{local-source} 中，向宿主给出真实路径和目标。例如：
\begin{code}
请使用cangjie-skill处理local-source/my-paper.pdf。
目标：服务新型布尔函数/多输出函数的设计。
先建立章节/构造/定理依赖图，重点拆构造思想、参数限制和证明步骤。
每条候选定位到页码与公式；一般证明缺口、开放问题及数表冲突要隔离。
保存日志，先交整书/整文骨架，不凭记忆补造原文。
\end{code}
仓颉大致经历：整书理解、五类提取、三重验证、独立入口晋级、RIA能力卡、关联与词典、压力测试、确定性编译\cite{cangjie}。用户在阶段门审阅的是具体结果，不是用确认代替来源验证。一次只做一个内容试点；中断后读取状态文件续跑，不必重复已经核查的材料。

要求每项构造输出四类证书：\textbf{输入契约}（域/基/位序/维数）、\textbf{生成结果}（真值表或确定性生成器）、\textbf{证明链}（局部谱到整体指标的每一步）、\textbf{计算范围}（哪些参数实测、哪些只属一般推导）。未被完整支持的内容进入待核库，不能为了增加卡数量而包装成active方法。

\section{读懂第一版的验证报告}
本版正常任务15/15完成、边界语义15/15，路由18/18；原严格字段断言28/30。两处error与not\_applicable状态名差异仍计严格不匹配，不能把“数学上正确拒绝”写成“所有字段断言通过”。单次对照也不是统计显著提升的证明。

第一版编号1.0.0是可追溯快照名，保留研究试用属性。尚未解决一般DC存在性、完整HML/FAI求解器、混合阵列一般精确界，也不保证所有密码攻击指标。源书并未随公共仓库分发，出处和来源哈希保留用于审计。

\section{把安装与研究过程写成日志}
建议每次实验保存如下记录，而不是只保存模型对话：
\begin{code}
version: 1.0.0
host: Codex or DeepSeek Harness
input: examples/mm.json
skill: mm-spectrum
domain_and_basis: explicit
command_and_exit_code: recorded
observed_metrics: from output JSON
proof_scope: instance / family / lower bound
unverified_items: explicit list
\end{code}
日志不能包含API key。输出文件内容与模型摘要冲突时，先复算程序和输入；不要为了让两者一致而手工改指标。修改算法后至少重跑相关正常例和边界例，再生成新版本。

\section{Codex 专用排错与项目使用}
\begin{longtable}{p{0.27\linewidth}p{0.64\linewidth}}
\toprule 现象 & 操作与成功标志 \\ \midrule
找不到技能 & 检查用户/项目扫描根，确保name与目录格式正确；重开会话，要求报告实际文件路径。\\
同名技能出现两次 & 查找用户级、项目级和旧兼容目录；只保留本次选定来源，先备份再移动。\\
不能运行Python & 在宿主终端核查解释器路径；使用已验证的Python绝对路径。系统终端可运行不代表宿主PATH相同。\\
只能看到文件但不能运行 & 检查当前沙箱/权限策略；先读与预演，依据宿主批准流程执行必要命令，不通过关闭所有限制解决。\\
模型说计算完成却无文件 & 要求保存输入JSON、输出JSON、命令、退出码和所有非零分量指标；人工打开输出核对。\\
\bottomrule
\end{longtable}
后续新论文或新函数候选，先用仓颉做来源与任务拆解，再用本项目做限定构造和验证。两者都要留日志；用同一个词“验证通过”时，务必说明是来源、公式推导、脚本实例还是在线宿主任务。
\section{升级、备份与跨平台操作}
\subsection{升级不覆盖旧研究}
先查看现有安装路径、版本和差异，再运行：
\begin{code}
python scripts/install.py --host codex --scope user --upgrade --dry-run
python scripts/install.py --host codex --scope user --upgrade
\end{code}
DeepSeek版本把host值改为 \file{deepseek}。若使用项目级，保留原来的scope和project参数。旧目录会存入扫描根的 \file{.backups/时间戳/技能名}，安装器报告具体位置；不自动删除历史研究目录。需要回退时先停止宿主，再将当前目录另存，恢复指定备份，重新验收。不要同时装两份同名旧版和新版来“试试”。

\subsection{macOS/Linux差异}
将 \file{python} 统一改为实际可用的 \file{python3}，安装和测试命令其他参数不变。下载仓颉可用：
\begin{code}
mkdir -p "$HOME/CryptoBoolFunc-Lab"
cd "$HOME/CryptoBoolFunc-Lab"
base=https://github.com/kangarooking/cangjie-skill/releases/download/v2.5.0
curl -fL "$base/cangjie-skill-2.5.0.zip" -o cangjie.zip
curl -fL "$base/cangjie-skill-2.5.0.zip.sha256" -o cangjie.sha256
shasum -a 256 cangjie.zip
cat cangjie.sha256
unzip cangjie.zip -d upstream
python3 -m pip install pyyaml
python3 upstream/cangjie-skill/scripts/cangjie.py doctor
\end{code}
逐字比较两个哈希后才继续。Linux可用 \file{sha256sum}。将完整 \file{upstream/cangjie-skill} 复制到所选宿主扫描根，或克隆本项目后用安装器的 \file{--cangjie-source} 参数一起安装。安装器通过Python解析用户目录，因此无需手工翻译Windows路径。

\section{重新编译本教程PDF}
安装带中文支持的TeX Live或MiKTeX。使用XeLaTeX，而不是pdfLaTeX。本教程采用Fandol中文字体与Latin Modern Mono，避免依赖私人电脑字体。
\begin{code}
python scripts/build_docs.py
\end{code}
此脚本对两份TeX各编译两遍，临时文件留在忽略的.build目录，PDF复制到docs。手动编译则先进入docs：
\begin{code}
xelatex -interaction=nonstopmode -halt-on-error tutorial-codex-v1.tex
xelatex -interaction=nonstopmode -halt-on-error tutorial-codex-v1.tex
\end{code}
另一份把文件名改为 \file{tutorial-deepseek-v1.tex}。成功不只看“生成PDF”：还要查看全部页的标题、中文、代码换行、链接和页码。日志中的Overfull、Missing character和LaTeX Error需要处理；不能把异常日志隐藏后继续发布。

\section{以后怎样按固定流程上传GitHub}
本仓库没有一份可以凭空套用的“统一最新强制上传规范”；项目采用核查日GitHub官方文档和自身规则建立流程\cite{githubpr,githubrelease}。后续按 \file{RELEASE_PROCESS.md} 执行：
\begin{enumerate}
\item 读取远端默认分支、最新SHA与保护规则，检查本地变更；先确定本次授权范围。
\item 更新版本、说明、数学边界及Codex贡献记录；只提交公开清单内的文件。
\item 运行测试，重编TeX/PDF并视觉验收；生成两个宿主ZIP、SHA256SUMS与manifest。
\item 建立 \file{codex/v版本-主题} 分支并提交；上传后创建PR，说明变化、验证与未测项。
\item 在授权和仓库规则允许时合并；有失败检查或审批要求则保持PR并说明条件，不强推。
\item 再读远端main/合并提交，核对PDF、ZIP和提交SHA，保存上传日志。
\end{enumerate}
原始PDF、全文抽取、个人绝对路径、密钥和编译缓存不进入公共提交。仓库的 \file{releases/v1.0.0} 是本项目版本资产目录；只有确实创建了GitHub Release和标签，才称创建了该平台的Release，二者不能混同。

\section{完成检查表}
\begin{enumerate}
\item 仓颉完整目录可读，doctor为PASS，版本和下载哈希有记录。
\item 当前宿主在正确工作区发现本项目入口，能读取15张卡和对应资源。
\item MM输出N=4；DS示例N=240且7分量；双码校正器N=448；D阵列固定表示得到10。
\item 非法输入被拒绝，DS一般R分支与其他待核项没有被重新启用。
\item 安装器、版本、PDF/TeX、SHA256与贡献者说明均可查；个人实验输出已另存。
\end{enumerate}

\section*{贡献与引用}
本项目由wernerchang发起和维护。Codex（OpenAI）作为AI协作贡献者，协助研究整理、程序、验证、教程和发布准备；具体记录见 \file{CONTRIBUTORS.md}。仓颉上游和《密码函数》作者分别署名。AI协作不替代维护者对学术结论的判断。

\begingroup\small
\begin{thebibliography}{9}
\setlength{\itemsep}{1pt}
\bibitem{cangjie} kangarooking. Cangjie Skill：\href{https://github.com/kangarooking/cangjie-skill}{官方仓库}及\href{https://github.com/kangarooking/cangjie-skill/releases/tag/v2.5.0}{v2.5.0刷新附件}。核查2026-11-04。
\bibitem{codexskills} OpenAI. \href{https://learn.chatgpt.com/docs/build-skills}{Build skills}。本地目录、完整技能与调用方式。
\bibitem{codexquick} OpenAI. \href{https://developers.openai.com/codex/quickstart}{Codex quickstart}。
\bibitem{dshreadme} DeepSeek AI. \href{https://github.com/deepseek-ai/deepseek-harness}{DeepSeek Harness官方README}，开发预览与启动方式。
\bibitem{dshskills} DeepSeek AI. \href{https://github.com/deepseek-ai/deepseek-harness/blob/master/packages/skill/skill-filesystem/README.md}{Filesystem skill provider}，扫描根及目录格式。
\bibitem{dshweb} DeepSeek AI. \href{https://deepseek-harness.github.io/deepseek-harness/en/guide/quickstart}{Web UI指南}，模型与工作区。
\bibitem{githubpr} GitHub Docs. \href{https://docs.github.com/en/pull-requests/how-tos/create-pull-requests/creating-a-pull-request}{Creating a pull request}。
\bibitem{githubrelease} GitHub Docs. \href{https://docs.github.com/en/repositories/releasing-projects-on-github/about-releases}{About releases}。
\end{thebibliography}
\endgroup
\end{document}

